% संपूर्ण मराठी ग्रंथातील प्रकरण 39: विन्यासमीमांसा आणि चिन्हार्थमीमांसा
% हा संपादनयोग्य स्रोतखंड आहे. संकलनासाठी संपूर्ण स्रोतसंग्रहातील
% अक्षररूपे, पूर्वभाग, चिन्हव्याख्या आणि आकृती-स्रोत आवश्यक आहेत.
% मूळ: Open Logic Project; मजकूर आणि रूपांतर CC BY 4.0.
% यंत्रानुवाद, दुरुस्ती आणि तपासणी: OpenAI Codex —
% GPT-5.6 Sol आणि GPT-6 Sol, दोन्ही Ultra effort.
% दुरुस्ती-पुनरावलोकन: GPT-6.1 Sol, Ultra effort.
\chapter{विन्यासमीमांसा आणि चिन्हार्थमीमांसा}\label{mvl:syn:chap}









\section{परिचय}\label{mvl:syn:int:sec}

अभिजात तर्कशास्त्रात $\lnot$, $\land$, $\lor$, $\lif$ आणि
$\liff$ ही विधानीय संयोजके वापरून विधानीय चले पासून
घडवलेल्या सूत्रांचा विचार करतो. अभिजात तर्कशास्त्राची
चिन्हार्थमीमांसा परिभाषित करताना आपण $\True$ आणि $\False$ ही दोन
सत्यतामूल्ये वापरतो. सत्य-मूल्यांकन~$\pAssign{v}$ मध्ये
विधानीय चले ना अर्थ देतो; ते विधानीय चले ना $\True$ आणि
$\False$ ही सत्यतामूल्ये नेमून देते. मग कोणतेही सत्य-मूल्यांकन
कोणत्याही सूत्र~$\varphi$ चे $\pValue{v}(\varphi)$ हे सत्यतामूल्य
ठरवते; आणि सूत्राची सत्य-मूल्यांकन~$\pAssign{v}$ मध्ये
पूर्ति होते, $\pSat{v}{\varphi}$, जर आणि फक्त जर
$\pValue{v}(\varphi) = \True$.

बहुमूल्य तर्कशास्त्रे $\True$ आणि $\False$ यांपलीकडील
सत्यतामूल्यांना परवानगी देऊन अभिजात द्विमूल्य तर्कशास्त्राचे
सामान्यीकरण करतात. म्हणून बहुमूल्य तर्कशास्त्रात
सत्य-मूल्यांकन~$\pAssign{v}$ हे प्रत्येक विधानीय चल~$p$
ला संभाव्य सत्यतामूल्यांपैकी एक मूल्य नेमून देणारे फलन असते.
अनुमत सत्यतामूल्यांच्या संचाला आपण साधारणपणे~$V$ म्हणू.
$V = \{\True, \False\}$ असलेले अभिजात तर्कशास्त्र हेही
बहुमूल्य तर्कशास्त्र आहे; त्यात $\pValue{v}(\varphi)$ हे सत्यतामूल्य
संयोजकांच्या परिचित वैशिष्ट्यपूर्ण सत्यताकोष्टकांनुसार ठरते.

अधिक सत्यतामूल्ये समाविष्ट केल्यावर ``आणि,'' ``किंवा,'' ``नाही,''
आणि ``जर---तर'' असे वाचत असलेल्या संयोजकांसाठी
$\pValue{v}(\varphi)$ कसे काढायचे, याचे एकाहून अधिक स्वाभाविक
पर्याय असतात. म्हणून बहुमूल्य तर्कशास्त्र केवळ सत्यतामूल्यांच्या
संचानेच नव्हे, तर प्रत्येक संयोजकासाठी निवडलेल्या
\emph{सत्यता-फलनांनी}ही निश्चित होते. तथापि, तर्कशास्त्र~$\Log L$
साठी ही निवड झाल्यावर, अभिजात तर्कशास्त्राप्रमाणेच
$\pValue{v}(\varphi)[\Log L]$ हे सत्यतामूल्य त्या सत्य-मूल्यांकनाने
एकमेव रीतीने ठरते. अभिजात तर्कशास्त्राप्रमाणे बहुमूल्य
तर्कशास्त्रेही \emph{सत्यता-फलनात्मक} आहेत.

हे चिन्हार्थविषयक पायाभूत घटक मिळाल्यावर सर्वतः सत्यता,
तार्किक निष्पन्नता आणि पूर्ततायोग्यता यांच्या समकक्ष
चिन्हार्थविषयक संकल्पना परिभाषित करता येतात. अभिजात
तर्कशास्त्रात सूत्र सर्वतः सत्य असते, जर
कोणत्याही~$\pAssign{v}$ साठी त्याचे सत्यतामूल्य
$\pValue{v}(\varphi) = \True$ असेल. बहुमूल्य तर्कशास्त्रात ही
संकल्पना थोडी व्यापक करावी लागते. प्रथम, सत्यतामूल्यांच्या~$V$
संचात $\True$ असणे आवश्यक नाही. उदाहरणार्थ,
काही बहुमूल्य तर्कशास्त्रे $0$ आणि~$1$ यांदरम्यानच्या सर्व
परिमेय संख्या सत्यतामूल्ये म्हणून वापरतात. अशा वेळी $1$
सहसा $\True$ ची भूमिका बजावते. इतर तर्कशास्त्रांत
एकच नव्हे, तर अनेक सत्यतामूल्ये ती भूमिका बजावतात.
म्हणून प्रत्येक बहुमूल्य तर्कशास्त्रात \emph{निर्दिष्ट
सत्यतामूल्यां}चा संच~$V^+$ असावा अशी अट घालतो.
मग सूत्राची सत्य-मूल्यांकन~$\pAssign{v}$ मध्ये
पूर्ति होते, $\pSat{v}{\varphi}[\Log L]$, जर आणि फक्त जर
$\pValue{v}(\varphi)[\Log L] \in V^+$. सूत्र~$\varphi$
हे त्या तर्कशास्त्रातील सर्वतः सत्य सूत्र असते,
$\Entails[\Log L] \varphi$, जर आणि फक्त जर
कोणत्याही~$\pAssign{v}$ साठी $\pValue{v}(\varphi) \in V^+$.
शेवटी, सूत्रांच्या संचापासून $\varphi$ निष्पन्न होते,
$\Gamma \Entails[\Log L] \varphi$, जर~$\Gamma$ मधील सर्व
सूत्रांची पूर्ति करणारे प्रत्येक सत्य-मूल्यांकन~$\varphi$
चीही पूर्ति करत असेल.












\section{भाषा आणि संयोजके}\label{mvl:syn:con:sec}

अभिजात विधानीय तर्कशास्त्रात आणि इतर अनेक तर्कशास्त्रांत
\emph{विधानीय स्थिर} व \emph{संयोजक} यांचा संच वापरला जातो.
उदाहरणार्थ, पुढील चिन्हे आपण आदिम मानतो:
\begin{enumerate}
\item असत्यता दर्शवणारा विधानीय स्थिर~$\lfalse$.
\item सत्यता दर्शवणारा विधानीय स्थिर~$\ltrue$.
\item तार्किक संयोजके:
  \startycommalist
  \ycomma $\lnot$ (नकारण)
  \ycomma $\land$ (संधियोग)
  \ycomma $\lor$ (विकल्पयोग)
  \ycomma $\lif$ (शर्त)
  \ycomma $\liff$ (द्विशर्त)
\end{enumerate}


बहुमूल्य तर्कशास्त्रांतही हीच संयोजके वापरली जातात. तथापि,
उदाहरणार्थ संधियोगाच्या वेगवेगळ्या आवृत्त्या एकाच तर्कशास्त्रात
असणे अनेकदा उपयुक्त ठरते; त्या आवृत्त्या वेगळ्या दाखवण्यासाठी
भिन्न चिन्हे लागतात. काही बहुमूल्य तर्कशास्त्रांत अभिजात
तर्कशास्त्रात समतुल्य संयोजक नसलेली संयोजकेही असतात.
म्हणून आपण नेहमीपेक्षा थोडी अधिक सामान्य चौकट घेऊ.

\begin{defn}
  \emph{विधानीय भाषा} ही \emph{संयोजकां}च्या $\Lang L$
  या संचाने बनते. प्रत्येक संयोजक~$\star$ ची \emph{स्थानसंख्या}
  असते; स्थानसंख्या~$n$ असलेला संयोजक \emph{$n$-स्थानी}
  म्हणतात. स्थानसंख्या~$0$ असलेल्या संयोजकांना
  \emph{स्थिर} असेही म्हणतात; स्थानसंख्या~$1$ असलेल्या
  संयोजकांना \emph{एकस्थानी}, आणि स्थानसंख्या~$2$ असलेल्या
  संयोजकांना \emph{द्विस्थानी} म्हणतात.
\end{defn}

\begin{ex}
  विधानीय तर्कशास्त्राच्या मानक भाषेत~$\Lang L_0$ पुढील
  संयोजके (त्यांच्या स्थानसंख्यांसह) असतात:
  $\lfalse$~($0$),
  $\lnot$~($1$),
  $\land$~($2$),
  $\lor$~($2$),
  $\lif$~($2$). आपण विचारात घेणारी बहुतेक तर्कशास्त्रे
  ही भाषा वापरतात. काही तर्कशास्त्रे प्रस्थापित संकेतांनुसार
  काही संयोजकांसाठी वेगळी चिन्हे वापरतात. उदाहरणार्थ,
  गुणाकार तर्कशास्त्रात संधियोगाचे चिन्ह $\land$ ऐवजी
  अनेकदा~$\odot$ असते. कधी नवीन कारक जोडणेही सोयीचे
  ठरते; उदाहरणार्थ, $\triangle$ हा निश्चितता कारक
  ($1$-स्थानी).
\end{ex}












\section{सूत्र}\label{mvl:syn:fml:sec}

\begin{defn}[सूत्र]
\label{mvl:syn:fml:defn:formulas}
विधानीय भाषा~$\Lang L$ मधील \emph{सूत्रे} चा संच~$\Frm[L]$
पुढीलप्रमाणे विगमनाने परिभाषित करतो:
\begin{enumerate}
\item प्रत्येक विधानीय चल~$\Obj p_i$ हे आण्विक
  सूत्र असते.
\item $\Lang L$ मधील प्रत्येक $0$-स्थानी संयोजक (विधानीय स्थिर)
  हे आण्विक सूत्र असते.
\item $\star$ हा $\Lang L$ मधील सकारात्मक स्थानसंख्येचा
  $n$-स्थानी संयोजक असेल आणि
  $\varphi_1$, \dots, $\varphi_n$ ही सूत्रे असतील, तर
  $\star(\varphi_1, \dots, \varphi_n)$ हे सूत्र असते.
\item यांखेरीज काहीही सूत्र नाही.
\end{enumerate}
$\star$ $1$-स्थानी असेल, तर $\star(\varphi_1)$ हे अनेकदा फक्त
$\star \varphi_1$ असे लिहितात. $\star$ $2$-स्थानी असेल, तर
$\star(\varphi_1,\varphi_2)$ हे अनेकदा $(\varphi_1 \star \varphi_2)$ असे लिहितात.
\end{defn}

नेहमीप्रमाणे, सर्वांत बाहेरचे कंस आपण अनेकदा न सांगता वगळू.

\begin{ex}
  मानक भाषेत~$\Lang{L_0}$,
  $\Obj p_1 \lif (\Obj p_1 \land \lnot \Obj p_2)$ हे सूत्र आहे.
  गुणाकार तर्कशास्त्राच्या भाषेत ते याऐवजी
  $\Obj p_1 \lif (\Obj p_1 \odot \lnot \Obj p_2)$ असे लिहितात.
  भाषेत $1$-स्थानी $\triangle$ जोडल्यास,
  $\triangle (\Obj p_1 \land \Obj p_2) \lif (\triangle \Obj p_1 \land \triangle \Obj p_2)$
  यांसारखी सूत्रेही मिळतात.
\end{ex}












\section{मॅट्रिक्स}\label{mvl:syn:mat:sec}

बहुमूल्य तर्कशास्त्र त्याची भाषा, सत्यतामूल्यांचा~$V$ संच,
निर्दिष्ट सत्यतामूल्यांचा उपसंच आणि त्याच्या संयोजकांसाठीची
सत्यता-फलने यांवरून परिभाषित होते. या घटकांना एकत्रितपणे
\emph{मॅट्रिक्स} म्हणतात.

\begin{defn}[मॅट्रिक्स]
\label{mvl:syn:mat:defn:matrix}
तर्कशास्त्र~$\Log L$ साठीच्या \emph{मॅट्रिक्स} मध्ये पुढील घटक असतात:
\begin{enumerate}
\item भाषा~$\Lang L$ बनवणारा संयोजकांचा संच;
\item सत्यतामूल्यांचा $V \neq \emptyset$ संच;
\item निर्दिष्ट सत्यतामूल्यांचा $V^+ \subseteq V$ संच;
\item $\Lang L$ मधील प्रत्येक $n$-स्थानी संयोजक~$\star$ साठी
  सत्यता-फलन~$\tf{\star} : V^n \to V$. $n = 0$ असल्यास
  $\tf{\star}$ हे फक्त~$V$ मधील एक मूल्य असते.
\end{enumerate}
\end{defn}

\begin{ex}
अभिजात तर्कशास्त्र~\LogCL{} साठीच्या मॅट्रिक्समध्ये पुढील घटक असतात:
\begin{enumerate}
  \item $\lfalse$, $\lnot$, $\land$, $\lor$, $\lif$ असलेली
  मानक विधानीय भाषा~$\Lang L_0$.
  \item सत्यतामूल्यांचा संच $V = \{\True, \False\}$.
  \item $\True$ हेच एकमेव निर्दिष्ट मूल्य, म्हणजे $V^+ = \{\True\}$.
  \item $\lfalse$ साठी $\tf{\lfalse} = \False$. इतर सत्यता-फलने
  नेहमीच्या सत्यताकोष्टकांनी दिली आहेत (पहा
  \ref{mvl:syn:mat:fig:tf-CL}).
\end{enumerate}
\begin{figure}
  \begin{center}
      \begin{tabular}{c|c}
        $\tf{\lnot}$ & \\
        \hline
        $\True$ & $\False$ \\
        $\False$ & $\True$
      \end{tabular}
      \quad
      \begin{tabular}{c|cc}
        $\tf{\land}$ & $\True$ & $\False$ \\
        \hline
        $\True$ & $\True$ & $\False$ \\
        $\False$ & $\False$ & $\False$
      \end{tabular}
      \quad
      \begin{tabular}{c|cc}
        $\tf{\lor}$ & $\True$ & $\False$ \\
        \hline
        $\True$ & $\True$ & $\True$ \\
        $\False$ & $\True$ & $\False$
      \end{tabular}
      \quad
      \begin{tabular}{c|cc}
        $\tf{\lif}$ & $\True$ & $\False$ \\
        \hline
        $\True$ & $\True$ & $\False$ \\
        $\False$ & $\True$ & $\True$
      \end{tabular}
    \end{center}
    \caption{अभिजात तर्कशास्त्र~$\LogCL$ साठीची सत्यता-फलने.}
    \label{mvl:syn:mat:fig:tf-CL}
  \end{figure}
\end{ex}












\section{सत्य-मूल्यांकन आणि पूर्ति}\label{mvl:syn:val:sec}

\begin{defn}[सत्य-मूल्यांकने]
$V$ हा सत्यतामूल्यांचा संच असू द्या. $\Lang{L}$ साठी~$V$
मध्ये जाणारे \emph{सत्य-मूल्यांकन} म्हणजे भाषेतील
विधानीय चले ना~$V$ मधील घटक नेमून
देणारे फलन~$\pAssign{v}$ होय; म्हणजेच
$\pAssign{v} \colon \PVar \to V$.
\end{defn}

\begin{defn}\label{mvl:syn:val:defn:pValue}
  बहुमूल्य तर्कशास्त्र~$\Log L$ च्या सत्यतामूल्यांच्या~$V$
  संचात जाणारे सत्य-मूल्यांकन~$\pAssign{v}$ दिले असता,
  मूल्यन फलन $\pValue{v} \colon \Frm[L] \to V$
  पुढीलप्रमाणे विगमनाने परिभाषित करतो:
  \begin{enumerate}
    \item $\pValue{v}(\Obj p_n) = \pAssign{v}(\Obj p_n)$;
    \item $\star$ हा $0$-स्थानी संयोजक असेल, तर
    $\pValue{v}(\star)  = \tf{\star}[\Log L]$;
    \item $\star$ हा सकारात्मक स्थानसंख्येचा
    $n$-स्थानी संयोजक असेल, तर
    \[      \pValue{v}(\star(\varphi_1, \dots, \varphi_n)) = \tf{\star}[\Log L]
      (\pValue{v}(\varphi_1), \dots, \pValue{v}(\varphi_n)).
    \]
  \end{enumerate}
\end{defn}

\begin{defn}[पूर्ति]
\label{mvl:syn:val:defn:satisfaction} सूत्र~$\varphi$ ची
  सत्य-मूल्यांकन~$\pAssign{v}$ ने \emph{पूर्ति होते},
  $\pSat{v}{\varphi}[\Log L]$, जर आणि फक्त जर
  $\pValue{v}(\varphi)[\Log L] \in V^+$; येथे $V^+$ हा
  तर्कशास्त्र~$\Log L$ मधील निर्दिष्ट सत्यतामूल्यांचा संच आहे.

  ``$\pSat{v}{\varphi}[\Log L]$ नाही'' हे दर्शवण्यासाठी आपण
  $\pSat/{v}{\varphi}[\Log L]$ असे लिहितो. $\Gamma$ हा
  सूत्रांचा संच असेल, तर $\pSat{v}{\Gamma}[\Log L]$
  जर आणि फक्त जर प्रत्येक~$\varphi \in \Gamma$ साठी
  $\pSat{v}{\varphi}[\Log L]$.
\end{defn}












\section{चिन्हार्थविषयक संकल्पना}\label{mvl:syn:sem:sec}

बहुमूल्य तर्कशास्त्र~$\Log L$ एखाद्या मॅट्रिक्सने दिलेले
आहे असे समजू. मग~$\Log L$ साठी नेहमीच्या चिन्हार्थविषयक
संकल्पना परिभाषित करता येतात.

\begin{defn}
\begin{enumerate}
\item सूत्र~$\varphi$ \emph{पूर्ततायोग्य} असते, जर
  एखाद्या~$\pAssign{v}$ साठी $\pSat{v}{\varphi}$;
  कोणत्याही $\pAssign{v}$ साठी $\pSat{v}{\varphi}$ नसेल, तर ते
  \emph{अपूर्ततायोग्य} असते;
\item सूत्र~$\varphi$ हे \emph{सर्वतः सत्य सूत्र}
  असते, जर सर्व सत्य-मूल्यांकने~$v$ साठी $\pSat{v}{\varphi}$;
\item $\Gamma$ हा सूत्रांचा संच असेल, तर
  $\Gamma \Entails \varphi$ (``$\Gamma$ पासून $\varphi$
  निष्पन्न होते'') जर आणि फक्त जर
  $\pSat{v}{\Gamma}$ असलेल्या प्रत्येक
  सत्य-मूल्यांकन~$\pAssign{v}$ साठी $\pSat{v}{\varphi}$;
\item $\Gamma$ हा सूत्रांचा संच असेल, तर
  $\pSat{v}{\Gamma}$ असलेले सत्य-मूल्यांकन~$\pAssign{v}$
  अस्तित्वात असल्यास $\Gamma$ \emph{पूर्ततायोग्य} असतो;
  अन्यथा $\Gamma$ \emph{अपूर्ततायोग्य} असतो.
\end{enumerate}
\end{defn}

अभिजात तर्कशास्त्रात जशी काही तथ्ये या संकल्पनांसाठी
लागू होतात, तशीच ती येथेही लागू होतात:

\begin{prop}
\label{mvl:syn:sem:prop:semanticalfacts}
\begin{enumerate}
\item $\varphi$ हे सर्वतः सत्य सूत्र असते, जर आणि फक्त जर
  $\emptyset \Entails \varphi$;
\item $\Gamma$ पूर्ततायोग्य असेल, तर $\Gamma$ चा प्रत्येक सांत
  उपसंचही पूर्ततायोग्य असतो;
\item\label{mvl:syn:sem:def:monotonicity}
एकस्वनिकता: $\Gamma \subseteq \Delta$ आणि
  $\Gamma \Entails \varphi$ असल्यास $\Delta \Entails \varphi$ देखील;
\item\label{mvl:syn:sem:def:Cut}
संक्रमकता: $\Gamma \Entails \varphi$ आणि
  $\Delta \cup \{ \varphi\} \Entails \psi$ असल्यास
  $\Gamma \cup \Delta \Entails \psi$;
\end{enumerate}
\end{prop}

\begin{proof}
सराव.
\end{proof}

\begin{prob}
\ref{mvl:syn:sem:prop:semanticalfacts} सिद्ध करा.
\end{prob}

अभिजात तर्कशास्त्रात निष्पन्नता आणि सशर्त संयोजक यांचा संबंध
जोडता येतो. उदाहरणार्थ, \emph{विध्यात्मक अनुमान} वैध असते:
जर $\Gamma \Entails \varphi$ आणि $\Gamma \Entails \varphi \lif \psi$
असेल, तर $\Gamma \Entails \psi$. अभिजात तर्कशास्त्रात
$\Entails$ आणि $\lif$ यांच्यातील दुसरा महत्त्वाचा संबंध म्हणजे
चिन्हार्थविषयक निगमन प्रमेय: $\Gamma \Entails \varphi \lif \psi$
जर आणि फक्त जर $\Gamma \cup \{\varphi\} \Entails \psi$.
हे निष्कर्ष बहुमूल्य तर्कशास्त्रांत \emph{नेहमीच लागू होतात असे नाही}.
ते लागू होणे $\tf{\lif}$ या सत्यता-फलनावर अवलंबून असते.












\section[अभिजात तर्कशास्त्राची उपतर्कशास्त्रे]{अभिजात तर्कशास्त्र~$\LogCL$ ची उपतर्कशास्त्रे म्हणून बहुमूल्य तर्कशास्त्रे}\label{mvl:syn:sub:sec}

नेहमीची बहुमूल्य तर्कशास्त्रे अशा मॅट्रिक्सांनी परिभाषित होतात की
$\{\True, \False\}$ मधील आदानांवर त्यांच्या सत्यता-फलनांची मूल्ये
अभिजात सत्यता-फलनांशी जुळतात. नेमके म्हणजे, या तर्कशास्त्रांत
$x \in \{\True, \False\}$ असल्यास $\tf{\lnot}[\Log L](x) =
\tf{\lnot}[\LogCL](x)$; आणि $\star$ हे $\land$, $\lor$,
$\lif$ यांपैकी कोणतेही असेल व $x, y \in \{\True, \False\}$
असतील, तर $\tf{\star}[\Log L](x,y) =
\tf{\star}[\LogCL](x,y)$. दुसऱ्या शब्दांत,
$\lnot$, $\land$, $\lor$, $\lif$ यांची सत्यता-फलने
$\{\True,\False\}$ वर मर्यादित केली असता ती नेमकी
अभिजात सत्यता-फलनेच असतात.

\begin{prop}\label{mvl:syn:sub:prop:mvl-cl}
  बहुमूल्य तर्कशास्त्र~$\Log L$ च्या भाषेत
  $\lnot$, $\land$, $\lor$, $\lif$ ही संयोजके आहेत,
  $\True, \False \in V$ आहेत आणि त्याची सत्यता-फलने
  पुढील अटी पूर्ण करतात असे समजा:
  \begin{enumerate}
    \item\label{mvl:syn:sub:prop:not} $x = \True$ किंवा $x = \False$
    असल्यास $\tf{\lnot}[\Log L](x) = \tf{\lnot}[\LogCL](x)$;
    \item\label{mvl:syn:sub:prop:land} $\tf{\land}[\Log L](x,y) = \tf{\land}[\LogCL](x,y)$,
    \item\label{mvl:syn:sub:prop:lor} $\tf{\lor}[\Log L](x,y) = \tf{\lor}[\LogCL](x,y)$,
    \item\label{mvl:syn:sub:prop:lif} $\tf{\lif}[\Log L](x,y) = \tf{\lif}[\LogCL](x,y)$,
    जेव्हा $x, y \in \{\True, \False\}$.
  \end{enumerate}
  मग $V$ मध्ये जाणाऱ्या कोणत्याही सत्य-मूल्यांकन~$\pAssign v$
  साठी, प्रत्येक विधानीय चलाबाबत $\pAssign
  v(p) \in \{\True,\False\}$ असल्यास,
  या चार संयोजकांनी आणि विधानीय चलांनीच बनलेल्या प्रत्येक
  सूत्रासाठी $\pValue v[\Log L](\varphi) = \pValue
  v[\LogCL](\varphi)$.
\end{prop}

\begin{proof}
  $\varphi$ वरील विगमनाने.
  \begin{enumerate}
  \item $\varphi \ident p$ हे आण्विक असेल, तर
  $\pValue v[\Log L](\varphi) = \pAssign v(p) = \pValue
  v[\LogCL](\varphi)$.
  \item $\varphi \ident \lnot B$ असेल, तर
  \begin{align*}
    \pValue v[\Log L](\varphi) & = \tf{\lnot}[\Log L](\pValue v[\Log L](\psi)) & &\text{\ref{mvl:syn:val:defn:pValue} नुसार}\\
    & = \tf{\lnot}[\Log L](\pValue v[\LogCL](\psi)) && \text{विगमन-गृहीतकानुसार}\\
    & = \tf{\lnot}[\LogCL](\pValue v[\LogCL](\psi))
&&\text{\ref{mvl:syn:sub:prop:not} या गृहीतकानुसार,}\\
&&&\text{कारण $\pValue v[\LogCL](\psi) \in \{\True, \False\}$,}\\
& = \pValue v[\LogCL](\varphi) &&\text{\ref{mvl:syn:val:defn:pValue} नुसार}.
\end{align*}
\item $\varphi \ident (\psi \land \chi)$ असेल, तर
\begin{align*}
  \pValue v[\Log L](\varphi) & = \tf{\land}[\Log L](\pValue v[\Log L](\psi), \pValue v[\Log L](\chi)) & &\text{\ref{mvl:syn:val:defn:pValue} नुसार}\\
  & = \tf{\land}[\Log L](\pValue v[\LogCL](\psi),\pValue v[\LogCL](\chi)) && \text{विगमन-गृहीतकानुसार}\\
  & = \tf{\land}[\LogCL](\pValue v[\LogCL](\psi),\pValue v[\LogCL](\chi))
&&\text{\ref{mvl:syn:sub:prop:land} या गृहीतकानुसार,}\\
&&&\text{कारण $\pValue v[\LogCL](\psi),\pValue v[\LogCL](\chi) \in \{\True, \False\}$,}\\
& = \pValue v[\LogCL](\varphi) &&\text{\ref{mvl:syn:val:defn:pValue} नुसार}.
\end{align*}
\end{enumerate}
$\varphi \ident (\psi \lor \chi)$ आणि $\varphi \ident (\psi \lif \chi)$
ही प्रकरणेही यांसारखीच आहेत.
\end{proof}

\begin{cor}
  या समान चार-संयोजक खंडातील सूत्रेच विचारात घेतली असता,
  बहुमूल्य तर्कशास्त्राने \ref{mvl:syn:sub:prop:mvl-cl} च्या अटी पूर्ण केल्या,
  $\True \in V^+$ आणि $\False \notin V^+$ असेल, तर
  ${\Entails[\Log L]} \subseteq {\Entails[\LogCL]}$;
  म्हणजे $\Gamma \Entails[\Log L] \psi$ असल्यास
  $\Gamma \Entails[\LogCL] \psi$. विशेषतः $\Log L$ मधील
  या खंडातील प्रत्येक सर्वतः सत्य सूत्र अभिजात तर्कशास्त्रातही
  सर्वतः सत्य असते.
\end{cor}

\begin{proof}
  व्यत्यय-विधान सिद्ध करू. $\Gamma \Entails/[\LogCL] \psi$
  असे समजा. मग असे सत्य-मूल्यांकन~$\pAssign v\colon \PVar \to
  \{\True, \False\}$ अस्तित्वात असते की सर्व
  $\varphi \in \Gamma$ साठी $\pValue v[\LogCL](\varphi) = \True$
  आणि $\pValue v[\LogCL](\psi) = \False$.
  $\True, \False \in V$ असल्याने सत्य-मूल्यांकन~$\pAssign v$
  हे~$\Log L$ साठीही सत्य-मूल्यांकन आहे.
  \ref{mvl:syn:sub:prop:mvl-cl} नुसार सर्व $\varphi \in \Gamma$ साठी
  $\pValue v[\Log L](\varphi) = \True$ आणि
  $\pValue v[\Log L](\psi) = \False$.
  $\True \in V^+$ आणि $\False \notin V^+$ असल्याने
  $\pSat{v}{\Gamma}[\Log L]$ आणि
  $\pSat/{v}{\psi}[\Log L]$; म्हणजे
  $\Gamma \Entails/[\Log L] \psi$.
\end{proof}


